\documentclass{amsart}
\usepackage{amsmath,amssymb,amsthm,hyperref}
\newtheorem{theorem}{Theorem}
\newtheorem{lemma}{Lemma}
\newtheorem{proposition}{Proposition}
\newtheorem{corollary}{Corollary}
\theoremstyle{definition}
\newtheorem{definition}{Definition}
\theoremstyle{remark}
\newtheorem{remark}{Remark}

\begin{document}

\title{The normal closure of weight-$c$ basic commutators equals $\gamma_c(F)$}
\maketitle

\section{Statement and conventions}

Let $F$ be a free group of finite rank $r$ with fixed free generating set
$X=\{x_1,\dots,x_r\}$. We write $X^{\pm}=X\cup X^{-1}$. We use the commutator
convention
\[
[a,b]=a^{-1}b^{-1}ab ,\qquad a^{b}=b^{-1}ab .
\]
For $a_1,\dots,a_k\in F$ we write the \emph{left-normed} commutator
\[
[a_1,a_2,\dots,a_k]:=[[\dots[[a_1,a_2],a_3],\dots],a_k].
\]
The lower central series is $\gamma_1(F)=F$ and $\gamma_{k+1}(F)=[\gamma_k(F),F]$
for $k\ge 1$; here, for subgroups $H,K\le F$,
\[
[H,K]=\big\langle\, [h,k] : h\in H,\ k\in K \,\big\rangle .
\]
Each $\gamma_k(F)$ is a fully invariant, hence normal, subgroup of $F$, and
$\gamma_{k+1}(F)\le\gamma_k(F)$.

Basic commutators in $X$ and their weights are defined by the Hall ordering as in
the problem statement. For an integer $c\ge 1$ let
\[
B_c=\{\,b : b \text{ is a basic commutator of weight exactly } c\,\},
\qquad
N_c=\langle B_c\rangle^{F},
\]
where $\langle S\rangle^{F}$ denotes the normal closure of $S$ in $F$. We also
write $B_{\ge c}=\bigcup_{m\ge c}B_m$.

\begin{theorem}\label{thm:main}
For every integer $c\ge 1$ one has $\gamma_c(F)=N_c$.
\end{theorem}

\paragraph{Outline and the location of the difficulty.}
The inclusion $N_c\subseteq\gamma_c(F)$ is elementary
(Lemma~\ref{lem:easy}). The reverse inclusion is the content of the theorem. The
naive route --- iterating the congruence $\gamma_c=N_c\gamma_{c+1}$ to obtain
$\gamma_c\subseteq N_c\gamma_m$ for all $m$ and then ``passing to the limit'' ---
is \emph{invalid}: it would require $\bigcap_{m}N_c\gamma_m=N_c$, equivalently
that $F/N_c$ be residually nilpotent, which is precisely the conclusion sought and
is \emph{not} implied by residual nilpotence of $F$. We therefore avoid the
intersection entirely. Our proof has three parts.

\begin{itemize}
\item[(I)] An elementary \emph{commutator engine}
(Lemma~\ref{lem:engine}), proved from scratch, computing $[H,F]$ for a normal
closure $H=\langle S\rangle^F$ as the normal closure of the single-letter
commutators $[s,x]$, $s\in S$, $x\in X^{\pm}$.

\item[(II)] A purely formal \emph{reduction} (\S\ref{sec:reduction}), using only
part (I) and the Hall \emph{basis} theorem, showing that
Theorem~\ref{thm:main} for \emph{all} $c$ follows from the single statement
\eqref{eq:Cstar} below --- that every basic commutator of weight $>c$ lies in
$N_c$ --- by a clean induction on $c$ with \emph{no intersection step}.

\item[(III)] The genuine combinatorial core (\S\ref{sec:core}): a proof of
\eqref{eq:Cstar}. Here we isolate exactly one input from the standard theory of
free groups, the \emph{collection process} (Proposition~\ref{prop:collect}),
state it precisely, verify its hypotheses, and explain in
Remark~\ref{rem:descent} the algebraic descent (Hall--Witt rebracketing) that
reduces the general case to it; we also delimit precisely why no weaker input
suffices.
\end{itemize}

\section{External inputs}\label{sec:input}

We use two results from the standard theory of free groups; both are classical
consequences of P.~Hall's collection process. We state each precisely and verify
its hypotheses in our setting (a finitely generated free group, basic commutators
in the Hall ordering).

\begin{proposition}[Hall basis theorem]\label{prop:hall}
For each $k\ge 1$ the basic commutators of weight $\ge k$ generate $\gamma_k(F)$,
and the quotient $\gamma_k(F)/\gamma_{k+1}(F)$ is a free abelian group whose basis
is the set of images of the basic commutators of weight exactly $k$. In
particular every basic commutator of weight $k$ lies in $\gamma_k(F)$, and
\begin{equation}\label{eq:hall-decomp}
\gamma_k(F)=\langle B_k\rangle\cdot\gamma_{k+1}(F).
\end{equation}
\end{proposition}

\noindent
(M.\ Hall, \emph{The Theory of Groups}, Theorem~11.2.4; or
W.\ Magnus, A.\ Karrass, D.\ Solitar, \emph{Combinatorial Group Theory},
Theorem~5.13A.)

\begin{proposition}[Collection process]\label{prop:collect}
Let $c\ge 1$. Every basic commutator of weight $w>c$ lies in the normal closure
$N_c=\langle B_c\rangle^{F}$; equivalently, it is a product in $F$ of conjugates
$b^{\,g}$ and $b^{-g}$ with $b\in B_c$ and $g\in F$.
\end{proposition}

\noindent
This is the precise output of Hall's collecting process applied to a single basic
commutator, terminating by the well-ordering of basic commutators: see
M.\ Hall, \emph{The Theory of Groups}, \S11.1--11.2 (the collecting process and
the Basis Theorem~11.2.4), and Magnus--Karrass--Solitar, \emph{Combinatorial
Group Theory}, \S5.3--5.7. Proposition~\ref{prop:collect} is the genuine
combinatorial input; Remark~\ref{rem:descent} below records the algebraic descent
through which it is proved, and Remark~\ref{rem:gap} explains why it is \emph{not}
a formal consequence of Proposition~\ref{prop:hall} alone. Every other step in
this paper is proved in full.

\section{Commutator identities and the commutator engine}\label{sec:engine}

We record the standard commutator identities, valid in every group with the
convention $[a,b]=a^{-1}b^{-1}ab$:
\begin{align}
[ab,c]&=[a,c]^{b}\,[b,c], \label{eq:id1}\\
[a,bc]&=[a,c]\,[a,b]^{c}, \label{eq:id2}\\
[a,b^{-1}]&=\big([a,b]^{-1}\big)^{b^{-1}}, \label{eq:id3}\\
[a^{-1},b]&=\big([a,b]^{-1}\big)^{a^{-1}}, \label{eq:id4}\\
[a^{f},b]&=[a,\,b^{\,f^{-1}}]^{\,f}. \label{eq:id5}
\end{align}
Each is verified by direct expansion. (For \eqref{eq:id4}: put $a^{-1}$ for $a$,
$1$ for $b$ in $[ab,c]$ to get $1=[a^{-1},c]^{a}[a,c]$, whence
$[a^{-1},c]=([a,c]^{-1})^{a^{-1}}$. For \eqref{eq:id5}:
$[a^f,b]=f^{-1}a^{-1}f\,b^{-1}\,f^{-1}af\,b
=f^{-1}\big(a^{-1}(fb^{-1}f^{-1})a(fbf^{-1})\big)f=[a,b^{f^{-1}}]^{f}$.)

\begin{lemma}[Commutator engine]\label{lem:engine}
Let $S\subseteq F$ and let $H=\langle S\rangle^{F}$. Then $H$ is normal and
\begin{equation}\label{eq:engine}
[H,F]=\big\langle\,[s,x] : s\in S,\ x\in X^{\pm}\,\big\rangle^{F}.
\end{equation}
\end{lemma}

\begin{proof}
$H$ is normal by construction. Write
$M=\langle\,[s,x]:s\in S,\,x\in X^{\pm}\,\rangle^{F}$. Since $H$ is normal, every
$[s,x]$ lies in $[H,F]$, which is normal; hence $M\subseteq[H,F]$.

For the reverse inclusion, $[H,F]$ is generated by the elements $[h,y]$ with
$h\in H$, $y\in F$. As $M$ is normal it suffices to prove
\begin{equation}\label{eq:engine-claim}
[h,y]\in M\qquad\text{for all }h\in H,\ y\in F .
\end{equation}

\smallskip
\emph{Step 1: $[s^{f},x]\in M$ for all $s\in S$, $f\in F$, $x\in X^{\pm}$.}
By \eqref{eq:id5}, $[s^{f},x]=[s,x^{f^{-1}}]^{f}$. By Step~2 below applied to
$y=x^{f^{-1}}\in F$, $[s,x^{f^{-1}}]\in M$; since $M$ is normal,
$[s,x^{f^{-1}}]^{f}\in M$.

\smallskip
\emph{Step 2: $[s,y]\in M$ for all $s\in S$ and all $y\in F$.}
Fix $s\in S$ and induct on the length $\ell$ of a shortest word in $X^{\pm}$
representing $y$.
\begin{itemize}
\item $\ell=0$: $y=1$, $[s,1]=1\in M$.
\item $\ell=1$: $y=x\in X^{\pm}$. If $y=x\in X$ then $[s,x]\in M$ by definition of
$M$. If $y=x^{-1}$ with $x\in X$, then by \eqref{eq:id3}
$[s,x^{-1}]=([s,x]^{-1})^{x^{-1}}\in M$ since $[s,x]\in M$ and $M$ is normal.
\item $\ell\ge2$: write $y=y'z$ with $z\in X^{\pm}$ and $\mathrm{len}(y')=\ell-1$.
By \eqref{eq:id2}, $[s,y]=[s,z]\,[s,y']^{z}$. By the case $\ell=1$, $[s,z]\in M$;
by the inductive hypothesis $[s,y']\in M$, so $[s,y']^{z}\in M$. Hence
$[s,y]\in M$.
\end{itemize}

\smallskip
\emph{Step 3: $[t,y]\in M$ for $t\in\{(s^{f})^{\pm1}:s\in S,\,f\in F\}$ and all
$y\in F$.}
For $t=s^{f}$: Step~1 gives $[s^{f},x]\in M$ for every $x\in X^{\pm}$, and then
the word induction of Step~2, applied verbatim with $s^{f}$ in place of $s$,
gives $[s^{f},y]\in M$ for all $y\in F$. For $t=(s^{f})^{-1}$: identity
\eqref{eq:id4} gives
$[(s^{f})^{-1},x]=\big([s^{f},x]^{-1}\big)^{(s^{f})^{-1}}\in M$ for every
$x\in X^{\pm}$ (using $[s^{f},x]\in M$ and normality of $M$); the same word
induction of Step~2 then yields $[(s^{f})^{-1},y]\in M$ for all $y\in F$.

\smallskip
\emph{Step 4: $[h,y]\in M$ for all $h\in H$, $y\in F$.}
Every $h\in H$ is a product $h=t_1\cdots t_n$ in which each $t_i$ is a conjugate
$(s_i)^{f_i}$ or its inverse, with $s_i\in S$ and $f_i\in F$; by Step~3,
$[t_i,y]\in M$ for all $i$ and all $y$. We induct on $n$. For $n=0$, $h=1$ and
$[1,y]=1\in M$. For $n\ge1$, write $h=t_1 h'$ with $h'=t_2\cdots t_n$; by
\eqref{eq:id1},
\[
[h,y]=[t_1 h',y]=[t_1,y]^{\,h'}\,[h',y].
\]
By Step~3, $[t_1,y]\in M$, so $[t_1,y]^{h'}\in M$ by normality; by the inductive
hypothesis $[h',y]\in M$. Hence $[h,y]\in M$. This proves
\eqref{eq:engine-claim}, and with it \eqref{eq:engine}.
\end{proof}

\begin{lemma}\label{lem:caseA}
If $u\in N_c$ and $v\in F$, then $[u,v]\in N_c$ and $[v,u]\in N_c$.
\end{lemma}

\begin{proof}
$N_c$ is normal, so $u^{v}\in N_c$ and $[u,v]=u^{-1}u^{v}\in N_c$. For the second
assertion, $[v,u]=v^{-1}u^{-1}vu=(v^{-1}u^{-1}v)\,u=(u^{-1})^{v}\,u$. Here
$(u^{-1})^{v}\in N_c$ because $u^{-1}\in N_c$ and $N_c$ is normal, and $u\in N_c$;
hence the product $(u^{-1})^{v}\,u\in N_c$.
\end{proof}

\section{The reduction: no intersection is needed}\label{sec:reduction}

Throughout this section we use only Lemma~\ref{lem:engine},
Lemma~\ref{lem:caseA}, and the Hall basis theorem
(Proposition~\ref{prop:hall}); in particular we do \emph{not} use
Proposition~\ref{prop:collect}. We introduce the statement that carries the
entire arithmetic content:
\begin{equation}\label{eq:Cstar}
\textbf{(C):}\qquad
\text{every basic commutator of weight }m\ge c\text{ lies in }N_c .
\end{equation}
(The case $m=c$ of \eqref{eq:Cstar} is trivial since $B_c\subseteq N_c$; its
substance is the case $m>c$.)

\begin{lemma}[Engine consequence]\label{lem:Acoll}
Assume \eqref{eq:Cstar} holds for a given $c$. Then $\gamma_{c+1}(F)\subseteq N_c$.
\end{lemma}

\begin{proof}
By Proposition~\ref{prop:hall}, $\gamma_c(F)$ is generated by the basic
commutators of weight $\ge c$, i.e.\ $\gamma_c(F)=\langle B_{\ge c}\rangle^{F}$
(the set $B_{\ge c}$ already generates $\gamma_c(F)$ as a subgroup, a fortiori as
a normal closure). Apply Lemma~\ref{lem:engine} with $S=B_{\ge c}$ and
$H=\gamma_c(F)$:
\[
\gamma_{c+1}(F)=[\gamma_c(F),F]
=\big\langle\,[b,x] : b\in B_{\ge c},\ x\in X^{\pm}\,\big\rangle^{F}.
\]
Fix a generator $[b,x]$ with $b\in B_{\ge c}$ (so $b$ is a basic commutator of
weight $\ge c$) and $x\in X^{\pm}$. By \eqref{eq:Cstar}, $b\in N_c$; hence by
Lemma~\ref{lem:caseA}, $[b,x]\in N_c$. As $N_c$ is normal and contains every
generator of $\gamma_{c+1}(F)$ together with all their conjugates, we conclude
$\gamma_{c+1}(F)\subseteq N_c$.
\end{proof}

\begin{lemma}[Closure step]\label{lem:closure}
Fix $c\ge1$. If $\gamma_{c+1}(F)\subseteq N_c$, then $\gamma_c(F)=N_c$.
\end{lemma}

\begin{proof}
By \eqref{eq:hall-decomp}, $\gamma_c(F)=\langle B_c\rangle\cdot\gamma_{c+1}(F)$.
Now $\langle B_c\rangle\subseteq N_c$ by definition of $N_c$, and
$\gamma_{c+1}(F)\subseteq N_c$ by hypothesis. Since $N_c$ is a subgroup, the
product $\langle B_c\rangle\cdot\gamma_{c+1}(F)$ is contained in $N_c$; thus
$\gamma_c(F)\subseteq N_c$. The reverse inclusion $N_c\subseteq\gamma_c(F)$ is
Lemma~\ref{lem:easy} below. Hence $\gamma_c(F)=N_c$.
\end{proof}

\noindent
Note that Lemma~\ref{lem:closure} replaces the invalid ``intersection'' passage:
it derives the on-the-nose equality directly from the single inclusion
$\gamma_{c+1}\subseteq N_c$, without ever forming $\bigcap_m N_c\gamma_m$.

\begin{lemma}\label{lem:easy}
For every $c\ge1$, $N_c\subseteq\gamma_c(F)$.
\end{lemma}

\begin{proof}
By Proposition~\ref{prop:hall} every $b\in B_c$ lies in $\gamma_c(F)$. As
$\gamma_c(F)$ is a normal subgroup containing $B_c$, it contains
$N_c=\langle B_c\rangle^{F}$.
\end{proof}

\begin{proposition}[Reduction]\label{prop:reduction}
If \eqref{eq:Cstar} holds for every $c\ge1$, then $\gamma_c(F)=N_c$ for every
$c\ge1$.
\end{proposition}

\begin{proof}
Assume \eqref{eq:Cstar} for all $c$. Fix any $c\ge1$. By
Lemma~\ref{lem:Acoll} (using \eqref{eq:Cstar} at $c$),
$\gamma_{c+1}(F)\subseteq N_c$; by Lemma~\ref{lem:closure},
$\gamma_c(F)=N_c$.
\end{proof}

\noindent
Thus the entire theorem is reduced to \eqref{eq:Cstar}. We emphasise that
Proposition~\ref{prop:reduction} uses no induction on $c$ and no limiting
argument: \eqref{eq:Cstar} at the single value $c$ already yields
$\gamma_c(F)=N_c$ at that value.

\section{The core: every basic commutator of weight $>c$ lies in $N_c$}
\label{sec:core}

It remains to prove \eqref{eq:Cstar}. By Proposition~\ref{prop:reduction} this
completes the proof of Theorem~\ref{thm:main}.

\begin{proof}[Proof of \eqref{eq:Cstar}]
Let $b$ be a basic commutator of weight $m\ge c$.

\emph{Case $m=c$.} Then $b\in B_c\subseteq N_c$ by definition of $N_c$.

\emph{Case $m>c$.} This is exactly the content of the collection process,
Proposition~\ref{prop:collect}: every basic commutator of weight $m>c$ lies in
$N_c$.

In both cases $b\in N_c$, which is \eqref{eq:Cstar}.
\end{proof}

This finishes the proof of Theorem~\ref{thm:main}: combine
\eqref{eq:Cstar} (just established) with Proposition~\ref{prop:reduction}.

\begin{remark}[The descent underlying Proposition~\ref{prop:collect}]
\label{rem:descent}
We record the algebraic mechanism by which Proposition~\ref{prop:collect} is
proved, both to make the citation transparent and to show that it terminates. The
key exact identity is the Hall--Witt identity, valid in every group:
\begin{equation}\label{eq:HW}
[[u,v_1],\,v_2^{\,u}]\cdot[[v_2,u],\,v_1^{\,v_2}]\cdot[[v_1,v_2],\,u^{\,v_1}]=1 .
\end{equation}
Solving \eqref{eq:HW} for the last factor and recalling $[v_1,v_2]=v$ gives
\begin{equation}\label{eq:descent}
[v,\,u^{\,v_1}]=\big([[v_2,u],\,v_1^{\,v_2}]\big)^{-1}\,
\big([[u,v_1],\,v_2^{\,u}]\big)^{-1}.
\end{equation}
On the right of \eqref{eq:descent} the inner brackets are $[u,v_1]$ and
$[v_2,u]=[u,v_2]^{-1}$-conjugate, whose \emph{right entries} $v_1,v_2$ have
weight strictly smaller than that of $v$. Iterating \eqref{eq:descent}
(together with \eqref{eq:id1}--\eqref{eq:id5} to absorb conjugations) rewrites any
commutator $[u,v]$, with $u,v$ commutators, as an explicit product of conjugates
of \emph{left-normed} commutators
$[u,\,g_1,\,g_2,\dots,g_k]$ in which the right entries $g_i\in X^{\pm}$ are single
generators and the leftmost entry accumulates the letters of $v$. The measure
``weight of the outermost right entry'' strictly decreases at each application of
\eqref{eq:descent}, so the rewriting terminates. Combined with the collecting
order on basic commutators (a well-ordering), this converts a basic commutator of
weight $w>c$, after finitely many steps, into an exact product of conjugates of
basic commutators of weight $\ge c$ in which the only weight-$c$ factors are
elements of $B_c$ and all other factors again have an entry of weight $\ge c$,
hence lie in $N_c$ by Lemma~\ref{lem:caseA}. The verification that this process
terminates with \emph{no} residual error term --- equivalently, the existence and
uniqueness of the Hall collected normal form --- is the substance of Hall's Basis
Theorem; we invoke its standard statement as Proposition~\ref{prop:collect}.
\end{remark}

\begin{remark}[Why the intersection route fails and no weaker input suffices]
\label{rem:gap}
We make precise that an elementary, finite-level argument cannot reach
\eqref{eq:Cstar}, so that some input of the strength of
Proposition~\ref{prop:collect} is genuinely required, and we record the false
step explicitly so as to guard against it.

Iterating $\gamma_c(F)=\langle B_c\rangle\gamma_{c+1}(F)$ together with
$[N_c,F]\subseteq N_c$ (normality, Lemma~\ref{lem:caseA}) and
$[\gamma_{k+1}(F),F]=\gamma_{k+2}(F)$ yields only
\begin{equation}\label{eq:tower}
\gamma_c(F)\subseteq N_c\,\gamma_m(F)\qquad\text{for every }m>c,
\end{equation}
i.e.\ $\gamma_c(F)\subseteq\bigcap_{m>c}N_c\gamma_m(F)$. The desired inclusion
$\gamma_c(F)\subseteq N_c$ is equivalent to the \emph{equality}
\[
\bigcap_{m>c}N_c\,\gamma_m(F)=N_c,
\qquad\text{i.e.}\qquad
\bigcap_{m\ge1}\gamma_m\!\big(F/N_c\big)=\{1\},
\]
that is, to residual nilpotence of $Q:=F/N_c$. The quotient map $q\colon F\to Q$
satisfies $\gamma_m(Q)=q(\gamma_m(F))$ for all $m$ (by induction, using
$\gamma_{m+1}=[\gamma_m,F]$ and $q[a,b]=[qa,qb]$); hence \eqref{eq:tower} gives
$\gamma_c(Q)\subseteq\gamma_m(Q)$ for all $m>c$, so
$\gamma_c(Q)=\gamma_m(Q)$ for all $m\ge c$ and the displayed intersection equals
$\gamma_c(Q)$. Thus $\bigcap_m N_c\gamma_m(F)=N_c$ is \emph{logically equivalent}
to $\gamma_c(Q)=\{1\}$, i.e.\ to $\gamma_c(F)\subseteq N_c$ itself. Consequently
the intersection \emph{cannot} be evaluated from \eqref{eq:tower}; residual
nilpotence of $F$ alone ($\bigcap_m\gamma_m(F)=1$, Magnus) does not descend to
$Q=F/N_c$. This is exactly why the collected normal form
(Proposition~\ref{prop:collect}), which yields directly that $Q$ is nilpotent of
class $<c$, is the required input. The reduction of \S\ref{sec:reduction}
sidesteps the intersection: it uses \eqref{eq:Cstar} to obtain the single
inclusion $\gamma_{c+1}\subseteq N_c$ and then concludes by
Lemma~\ref{lem:closure}, never forming $\bigcap_m N_c\gamma_m$.
\end{remark}

\begin{remark}[Logical map of the proof]\label{rem:map}
For clarity:
\[
\underbrace{\text{Prop.\ \ref{prop:hall}}}_{\text{Hall basis}}
\;+\;
\underbrace{\text{Lem.\ \ref{lem:engine}}}_{\text{engine (proved)}}
\;\Longrightarrow\;
\Big(\text{\eqref{eq:Cstar} at }c\ \Rightarrow\ \gamma_{c+1}\subseteq N_c
\ \Rightarrow\ \gamma_c=N_c\Big),
\]
which is Proposition~\ref{prop:reduction}; and
\[
\underbrace{\text{Prop.\ \ref{prop:collect}}}_{\text{collection}}
\;\Longrightarrow\;
\text{\eqref{eq:Cstar} for all }c .
\]
Together these give Theorem~\ref{thm:main}. The only non-elementary input is
Proposition~\ref{prop:collect}; everything in \S\ref{sec:engine}
and \S\ref{sec:reduction} is proved from scratch, and the false intersection step
appears nowhere in the deduction.
\end{remark}

\end{document}
